\documentclass[a4paper]{article}
\usepackage{luatexja}
\usepackage{amsmath}
\usepackage{amssymb}
\usepackage{geometry}
\usepackage{hyperref}

\geometry{margin=25mm}

\title{符号つき静電力から導かれる電子軌道：\\
共有結合・金属結合・結晶}

\author{https://github.com/hypernumbernet}
\date{\today}

\begin{document}
\maketitle

\begin{abstract}
等スケール軌跡の上で、二重時空理論の電磁力はクーロン力を干渉因子 \(\gamma_s\) で割ったものである。\(\gamma_s\) の符号が斥力である。異符号の対は \(\gamma_s>0\) で引き合い、\(\gamma_s<0\) で斥け合う。同符号の対は逆である。引力を打ち消すために遠心力の項は加えない。電子と陽子の一対は、両方のチャネルを担う。二つの干渉因子の符号が同じところでは、二つの質量による重力の項と電気の項は同じ向きを向き、動径方向の力は消えない。ボーア半径では両方の因子が正であり、質量の項は電荷の項より、\(4.40\times 10^{-40}\) と \(4.42\times 10^{-40}\) のあいだに厳密にある因子だけ小さい。外側の井戸では、円軌道は \(5\ell/6\) と \(\ell\) のあいだに厳密にある半径より外で動径方向に戻り、その半径と第一の節のあいだでは不安定である。その井戸の円軌道のエネルギーは、同じ半径で最も小さく、値は負である。正の核電荷をもつ核のまわりの電子一つは、その半径を保ち、軌道のエネルギーは陽子一つの値にその電荷を掛けたものである。直径の両端に置いた二電子、正四面体の頂点に置いた四電子、正八面体の頂点に置いた六電子、立方体の頂点に置いた八電子は、第一の斥力殻に動径の釣り合いをちょうど一つずつもち、直径は角の極大である。\(2\)、\(8\)、\(18\)、\(32\) は次数 \(0\)、\(1\)、\(2\)、\(3\) までのモードに、双対スピノルとそのスピンを反転した相手の席を一つずつ置いた数であり、節の番号ではない。次数ごとに見ると、席は \(2\)、\(6\)、\(10\)、\(14\) である。席は方向の上の模様である。電子軌道と呼ばれるものは、軌道ではなく電子配置である。電子が動くことはあり、共有結合では位置が固定される。単極だけが方向によらず、それより高い席は、ある方向では正、別の方向では負である。直径、正四面体、正八面体、立方体の、半径を戻す点は \(\gamma_s<0\) にある。外側の井戸では、同じ二電子は、核電荷が少なくとも \(1/4\) ならば、どのような配置でも同時には力を消さない。向かい合う射線の上では、斥力は有限な半径のすべてで単位核電荷の引力の四分の一より厳密に小さい。間隔の干渉因子が正であるところでは、対蹠点は角の極小であり、電子は広がり、半径は戻らない。核が引くところでの二電子の殻は、単極の席が二つであることであり、符号つきの力の零ではない。外側の井戸では核が引く。正四面体、正八面体、立方体、正二十面体は、核電荷がそれぞれ少なくとも \(1\)、\(2\)、\(3\)、\(5\) ならば動径の根をもたない。どの弦も半径より長く、外向き係数はその電荷を下回る。正十二面体は半径より短い弦をもち、二十の頂点の遠方での係数は \(15/2\) と \(8\) のあいだに厳密にある。その弦が第一の節に達するより前、応答は \(7\) を厳密に上回り、核電荷が \(7\) 以下ならばその区間の全体で外へ押される。核電荷 \(Z\ge 8\) では、応答がその電荷に等しくなる位相が、その弦が第一の節に達するより前に、ちょうど一つある。そこでの半径は \(7\ell/10\) と \(10\ell/7\) のあいだに厳密にある。二十は頂点数であり、四つの主殻の数の一つではない。第一の節の半径の二倍より長い線分の上では、両方の核に対して外側の井戸にある電子は、中点からより大きい電荷の方へ引かれる。打ち消しは一つであり、より小さい電荷に近く、戻されない。陽子が二つ、電子が二つのときには、第一の斥力殻の中のちょうど一つの正方形で消え、その正方形は対角線を戻す。陽子の間隔と電子の間隔がともに外側の井戸にあるあいだは、その二つの間隔のどの比でも、二つの力は同時には消えない。電子と陽子の間隔が外側の井戸にあるあいだも、二つの力は同時には消えない。ヘリウムは電荷二の核と二つの電子であって二電子の殻のままであり、各核に電荷二を置いた同じ族は外側の井戸に平衡をもたず、二つの干渉因子が零でないかぎり正方形も平衡ではない。水分子は、核を原点に、二本の射線のそれぞれに、陽子より核に近い結合電子を一つ置く。それらの電子が斥け合うあいだ、接線方向の力は、さらに加わる電子のすべてが射線の上に立つか、もう一方の結合に向いた側から斥けるかぎり、より広い角の側に留まる。直線は、それらの電子を戻す。直線の上では、各電子が核とその陽子のあいだにあり、核電荷が少なくとも \(1/4\) ならば、外向きの力の和は負である。電子が核から陽子の少なくとも四分の一の距離にあり、陽子どうしが各陽子ともう一方の結合電子と同じだけ離れる間隔がなお外側の井戸にあるとき、各陽子に働く接線方向の力は、その等間隔と直線のあいだの角で消える。その等間隔における射線のあいだの角は、直線の三分の一より厳密に小さい。正四面体の射線の上の四つの共有電子対は、電子が一つの距離に、陽子がそれより大きい一つの距離にあるとき、接線方向の力を及ぼさない。その角の余弦は \(-1/3\) であり、角は直角より大きい。非共有電子対の射線には陽子がなく、その電子は内へ押される。その電子が隣の核からの距離の少なくとも \(1/\sqrt{2}\) 倍に留まり、間隔が外側の井戸にあるあいだ、結合距離にある共有電子対より強く結合を押しのけ、そのような対が二つあるときの狭める寄与は一つの二倍である。陽子は一方の対角線上に、電子は他方の対角線上に座る。そこでは静電ポテンシャルのヘッシアンが正定値である。同符号の電荷は対角線に沿って引き合い、異符号の電荷は辺に沿って斥け合う。対角線は辺の \(\sqrt{2}\) 倍であり、対角線の力は \(4.3532\,k e^{2}/\ell^{2}\)、対角線の剛性は \(8.178\,k e^{2}/\ell^{3}\) である。その対の分数 \(2/z\) だけを運ぶ最近接結合には、静電力の各成分が消える点がさらにある。近傍が八個および十二個のとき、それらの点は鞍点であり、部分力は共有結合の対角線に対して比 \(0.04323\) と \(0.1704\) にある。どのイオンに対しても外側の井戸にあるところでは、イオンだけでは電子を保持しない。互いに垂直な三つの方向に電子を変位させると、それぞれの方向に沿う力が変わり、変位をさらに押し進める三つの割合の和は正である。クーロンの法則ではその和は零である。さらに加わる電子は、それぞれが外側の井戸にあって別々のイオンより遠くにあるかぎり、これを覆さない。正八面体、立方体、立方八面体に並んだイオンの中心では、電子に働く力は消え、どの変位もさらに押し進められる。同じ殻は、正のイオンのどの変位も戻す。イオンは位置に保持され、電子は自由である。同符号のイオンの単純立方、体心立方、面心立方の、立方体で切った部分は、どの位相も \(\pi\) より小さく第一の節の上にないかぎり、殻をいくつ数えても、イオンをすべての方向に保持し、電子を押し出す。割合は \(\sin x\) の符号をもち、半径 \(\ell/(2\pi)\) で反転する。岩塩型、塩化セシウム型、閃亜鉛鉱型、および蛍石型の両方の格子点では、最も近い二つの殻が、ありふれたイオン電荷に対して、外側の井戸の全体で正のイオンを押し出す。\(2\pi/n\) の回転が三次元の格子を保つのは \(n\in\{1,2,3,4,6\}\) に限り、結晶の格子点は五回軸をもたない。その正方形の辺はボーア半径ではない。その辺をその半径に置くと、第一の節は \(91/40\) と \(12/5\) ボーア半径のあいだに厳密に来る。ボーア半径での実験室の符号は引力であるから、その半径は外側の井戸にあり、正方形は反転した層にある。同じ層状の符号を重力チャネルで読むと、それは共有された席の排他である。二つの電子質量を第一の重力節の内側に置くと、それらは単純極に出会う。\(G m_e^{2}\) の小ささはその近傍を薄くし、極は消さない。半径を戻す点は外向きの力の符号の境界であり、電子はそこを横切って振動する。正方形の上では、電子の波長と対角線の波長との比は \(45.08\) である。電子の座標における \(K/m\) と対角線における \(K/m\) との比の平方根であり、辺の同定をこの比は要しない。遷移の振動数は固有値の差であり、温度に比例して伸び縮みできない。定常軌道は電子の運動を固定する。陽子に対する電子の速度の差は、軌道が計算する量ではなく、温度とともに変わらない。調和振動子の梯子の上では、基底状態より上の量子の平均個数は \(1/(e^{x}-1)\) である。ここで \(x\) は \(k_B T\) を単位とした間隔である。大きな空洞では、モードの数が振動数の二乗に比例して増える。この増え方は空洞の幾何である。波長の密度は、そのとき \(k_B T\) を単位として \(9/2\) と \(5\) のあいだにただ一つの最大をもつ。実験室の値 \(4.96\,k_B T\) がその最大であり、したがってピークの波長と温度との積は定数である。積分した振動数密度は \(T^{4}\) に比例し、一つの古典的な根の電力は \(T\) に比例する。定数 \(h\)、\(k_B\)、\(c\)、およびモード数の増え方は、入力である。
\end{abstract}

\section{はじめに}

二重時空理論の電子殻の節は、分母にスカラー干渉因子 \(\gamma_s\) をもつ静電力を書く~\cite{dst-main}。大きな間隔では因子は \(1\) に近づき、クーロンの法則が回復される。より近くでは、\(\gamma_s\) は離散的な半径の集合で零を通り、それらの節のあいだで負になる。節は力の発散である。節の外側での動径ポテンシャルは下に有界でない。

純粋に引力的なクーロン力は、有限な半径では消えない。古典的な円軌道は、それを遠心項 \(L^{2}/(m r^{3})\) で打ち消す。保存する角運動量の運動学である。その項はここでは加えない。第\ref{sec:nocent}節は、外側の井戸の円軌道が動径方向に戻る側から不安定な側へ移る半径を、\(5\ell/6\) と \(\ell\) のあいだに厳密に置く。その半径は、それらの円軌道のエネルギーが最も小さい場所であり、エネルギーは負である。正の核電荷をもつ核のまわりの電子一つでは、半径は同じであり、エネルギーはその電荷を陽子一つの値に掛けたものである。外向きの寄与は、符号つきの力のうちにすでにある。\(\gamma_s\) が負ならば異符号の対は斥け合い、\(\gamma_s\) が正ならば同符号の対は斥け合う。これらの寄与の打ち消しは、電荷と幾何がそれを与えるときの動径の根である。主殻の収容数は第\ref{sec:shells}節の席であり、二電子の殻はそのような根ではない。クーロン引力と遠心障壁との釣り合いでもない。

本稿はその軌跡を、陽子二つと電子二つの二中心配置について計算する。外側の井戸では、その配置は間隔のどの比でも釣り合わない。電子と陽子の間隔が外側の井戸にあるあいだも、二つの力は同時には消えない。第一の斥力殻にある一つの正方形は、対角線を戻す。ヘリウムは電荷二の核と二つの電子であって二電子の殻のままであり、各核に電荷二を置いた同じ族は外側の井戸に平衡をもたず、二つの干渉因子が零でないかぎり正方形も平衡ではない。第\ref{sec:water}節は、核から出る二本の射線のそれぞれに、陽子より核に近い結合電子を一つ置く。それらの電子が斥け合うあいだ、接線方向の力は、さらに加わる電子のすべてが射線の上に立つか、もう一方の結合に向いた側から斥けるかぎり、より広い角の側に留まり、直線はそれらの電子を戻す。直線の上では核電荷が少なくとも \(1/4\) ならば外向きの力の和は負である。電子が核から陽子の少なくとも四分の一の距離にあり、陽子の等間隔がなお外側の井戸にあるとき、各陽子に働く接線方向の力は、その等間隔と直線のあいだで消え、等間隔における射線のあいだの角は直線の三分の一より厳密に小さい。第\ref{sec:angle}節は、正四面体の射線の上に四つの電子対を置く。その角の余弦は \(-1/3\) であり、角は直角より大きい。電子が一つの距離に、陽子がそれより大きい一つの距離にある四つの共有電子対は、接線方向の力を及ぼさない。非共有電子対の射線には陽子がなく、その電子は内へ押される。その電子が隣の核からの距離の少なくとも \(1/\sqrt{2}\) 倍に留まるあいだ、共有電子対より強く残る結合を押しのけ、そのような電子が二つあるときの狭める寄与は一つの二倍である。最近接結合が電子対の \(2/z\) だけを運ぶ場合も計算する。第\ref{sec:crystal}節は、同じ割合を結晶の格子点の上で読む。第\ref{sec:free}節は、どのイオンに対しても外側の井戸にあるところではイオンだけでは電子を保持せず、同符号のイオンの対称な殻はその中心にイオンを保持することを示す。金属の電子は自由である。第\ref{sec:crystal}節は、その保持を立方格子の打ち切りのすべての殻に広げ、ありふれたイオン結晶の最も近い二つの殻がイオンを押し出すことを示す。もう一つの安定性は重力チャネルに属する。本文はパウリ原理を、そのチャネルの第一斥力層として読む。二つのフェルミオンが重なるときに生成されるのであって、別の統計的要請ではない。電子に適用すると、同じ層が二つの質量を一つの軌道の席からはじき出す。電子質量は小さく、はじき出す力は極である。等スケールの同定 \(\alpha=\beta=\ell/r\) は、核の層とクーロン層にすでに用いられている読みである。電磁の長さは、電子の換算コンプトン波長である。ボーアの \(n^{2}\) 則も微細構造定数も導出しない。殻の数も同じである。\(2\)、\(8\)、\(18\)、\(32\) は単極、双極、四重極、八重極に双対スピノルとそのスピンを反転した相手の席を一つずつ置いた数である。単極だけがどの方向でも同じである。それより高い席は、ある方向では正、別の方向では負である。電子軌道と呼ばれるものは、軌道ではなく電子配置である。電子が動くことはあり、共有結合では位置が固定される。量子力学も電子軌道を定常な模様としてすでに読み、符号つきの力はその読みを位置そのものまで運ぶ。一中心の直径、正四面体、正八面体、立方体が閉じるのは半径だけであり、しかも反転した層の中だけである。外側の井戸にある二電子は、核電荷が少なくとも \(1/4\) ならば、同時には力を消さない。核が引く外側の井戸は、正四面体、正八面体、立方体、正二十面体に、核電荷がそれぞれ少なくとも \(1\)、\(2\)、\(3\)、\(5\) であるとき、動径の根を与えない。二十の頂点をもつ正十二面体は、核電荷が少なくとも \(8\) のとき、外側の井戸に動径の根をちょうど一つもち、その半径は \(7\ell/10\) と \(10\ell/7\) のあいだに厳密にある。二十は四つの主殻の数の一つではない。その遠方での係数は \(15/2\) と \(8\) のあいだに厳密にあり、最も近い弦が第一の節に達するより前、応答は \(7\) を厳密に上回る。核電荷が \(7\) 以下ならば、その区間の全体で外へ押される。第\ref{sec:eneg}節は、その井戸で二核のあいだの線分に電子を置く。中点からはより大きい電荷の方へ引かれ、打ち消しは一つで、より小さい電荷に近い。半径を戻す点は、同時に符号の境界でもある。第\ref{sec:heat}節は、それを横切る動径の振動を読む。振動数は \(\omega^{2}=K k e^{2}/(m\ell^{3})\) である。符号つきの釣り合いは一つであり、動く座標ごとに慣性が分かれる。遷移の振動数は固有値の差であり、温度を含まない。定常軌道は電子の運動を固定する。陽子に対する電子の速度の差は、軌道が計算する量ではなく、温度とともに変わらない。連続スペクトルは、空洞がモードの数を与えたときの、間隔の固定された梯子の占有である。実験室の数値が入るところは、入力であると記す。

\section{符号つきの力}
\label{sec:force}

一軸の上で、双対ローターの干渉因子は
\[
\gamma_s
  =\cosh\frac{\alpha}{2}\cos\frac{\beta}{2}
  -\sinh\frac{\alpha}{2}\sin\frac{\beta}{2}
\]
である。等スケール軌跡は \(\alpha=\beta=\ell/r\) と置く。無次元変数 \(x=\ell/(2r)\) では
\begin{equation}
\label{eq:gamma}
\gamma_s(x)=\cosh x\cos x-\sinh x\sin x,
\qquad
\gamma_s'(x)=-2\cosh x\sin x
\end{equation}
である。したがって \(\gamma_s(0)=1\) であり、各区間 \([n\pi,(n+1)\pi]\) にちょうど一つの単純零点がある。最初の零点 \(x_1\) は
\begin{equation}
\label{eq:node}
\tan x=\coth x
\end{equation}
の \((\pi/4,\pi/2)\) における唯一の根である。数値的には \(x_1=0.93755\) であるから、最外の節は
\[
\frac{r_1}{\ell}=\frac{1}{2x_1}=0.53330
\]
にある。第二の零点は \(x_2=3.9274\) にあり、第一の斥力殻の内側の縁は \(r_2/\ell=0.12731\) である。半径の大きい方から読むと、外側の井戸 \(r>r_1\) では \(\gamma_s>0\)、殻 \(r_2<r<r_1\) では \(\gamma_s<0\) であり、符号はこのあと交替する。第一の殻の内部での \(\gamma_s\) の極値は \(\gamma_s(\pi)=-\cosh\pi\) であり、\(r=\ell/(2\pi)\) にある。力の極値は、より大きい半径にあり、第\ref{sec:nocent}節に記す。

二電荷の間隔を増す向きの力は
\begin{equation}
\label{eq:force}
F_r=\frac{k q_1 q_2}{\gamma_s(r)\,r^{2}},
\qquad
k=\frac{1}{4\pi\epsilon_0}
\end{equation}
である。同符号の電荷は \(\gamma_s>0\) で斥け合い、\(\gamma_s<0\) で引き合う。異符号では逆である。どちらの向きも同じ式である。斥力を生むために、これ以上のものは加えない。

電子と陽子は、この電気の項の隣に、二つの質量による重力の項を保つ。\(\gamma_{s,G}\) を、あらゆる質量の重力を支配する長さ、荷電パイオンのコンプトン波長 \(\lambda_\pi\) における質量チャネルの干渉因子とすると、間隔を増す力は
\begin{equation}
\label{eq:both}
F_r
  =-\frac{G m_e m_p}{\gamma_{s,G}(r)\,r^{2}}
  -\frac{k e^{2}}{\gamma_s(r)\,r^{2}}
\end{equation}
である。二つの因子の符号が同じところでは、二つの項は同じ向きを向き、和は消えない。ボーア半径は両方の長さの外側の平坦部にある。そこでは両方の因子が正であり、両方の項が引き合い、質量の項は電荷の項より、\(4.40\times 10^{-40}\) と \(4.42\times 10^{-40}\) のあいだに厳密にある因子だけ小さい。

原点でのテイラー展開は \(\gamma_s''(0)=-2\) と \(\gamma_s''''(0)=-4\) から従い、
\begin{equation}
\label{eq:series}
\gamma_s(x)=1-x^{2}-\frac16 x^{4}+O(x^{6})
\end{equation}
である。ゆえに
\[
\frac{1}{\gamma_s(r)\,r^{2}}
  =r^{-2}+\frac{\ell^{2}}{4}\,r^{-4}+\frac{7\ell^{4}}{96}\,r^{-6}+O(r^{-8}).
\]
クーロン力への補正は \((\ell/r)^{2}\) の次数から始まる。

\(r>r_1\) では、\eqref{eq:force} を積分し無限遠で消える対ポテンシャルが得られる。置換 \(s=\ell/(2x)\) は \(ds/s^{2}=-(2/\ell)\,dx\) を与えるので、
\begin{equation}
\label{eq:potential}
\int_r^{\infty}\frac{ds}{\gamma_s(s)\,s^{2}}
  =\frac{2}{\ell}\int_0^{\ell/(2r)}\frac{dx}{\gamma_s(x)}
\end{equation}
である。ただし積分路は外側の井戸に留まる。上端が \(x_1\) に近づくと積分は発散する。殻の内側へ同じ式は続かない。無限遠からの路が極を横切るからである。殻の内部での局所的な力の釣り合いは、なお定義される。\eqref{eq:force} が要求するのは \(\gamma_s\neq 0\) だけだからである。

\section{遠心項を加えない}
\label{sec:nocent}

陽子が一つ、電子が一つの場合を考える。動径方向の力は
\[
F_r=-\frac{k e^{2}}{\gamma_s(r)\,r^{2}}
\]
である。\(\gamma_s\) が有限で零でないところでは、\(F_r\) は零でない。外側の井戸では、あらゆる有限半径で引力である。第一の斥力殻では、あらゆる有限半径で斥力である。静電力が消える半径はなく、それを手で導入することもない。

主論文の向心的な書き換え \(v^{2}=k e^{2}/(m\gamma_s r)\) は、外側の井戸で保存する角運動量が何を寄与するかを記録する。円形のニュートン軌道に対する運動学的な恒等式であり、\eqref{eq:force} の中の余分な力ではない。\(\gamma_s>0\) のとき、内向きの力はその加速度を担う。

引力を正に数えると、無次元の力は \(Y=4x^{2}/\gamma_s(x)\) である。外側の井戸ではこれは正であり、位相について厳密に増える。半径が下がるにつれて引力は強くなる。井戸の境界は節である。第一の斥力殻の内部では、同じ力は負である。どちらの節でも発散し、極値は一つだけある。その位相は \(5/2\) と \(\pi\) のあいだに厳密にあり、半径は \(\ell/(2\pi)\) と \(\ell/5\) のあいだに厳密にある。そこでの外向きの大きさは \(3\,k e^{2}/\ell^{2}\) より厳密に小さい。\(x=\pi\)、すなわち \(\gamma_s=-\cosh\pi\) では、
\begin{equation}
\label{eq:plateau}
\bigl|F\bigr|
  =\frac{4\pi^{2}}{\cosh\pi}\,\frac{k e^{2}}{\ell^{2}}
  =3.4057\,\frac{k e^{2}}{\ell^{2}}
\end{equation}
であり、向きは外向きである。\(1/r^{2}\) の増大は、より弱い極値を、この点より大きい半径へ置く。角運動量がなければ、異符号の対は、第一の節のどちら側からもその節へ向かう。外側の井戸は内へ引き、第一の殻は外へ押す。同符号の対はその節から離される。実験室の井戸で殻を閉じる壁は、その斥力である。一中心の電子は、力の消える円の上には座らない。力の消える軌跡が現れるのは、複数の電荷があり、それらの符号つきの寄与が打ち消し合うときである。外側の井戸で二核のあいだに現れる打ち消しは席ではない（第\ref{sec:eneg}節）。

円軌道が動径方向に戻るのは、対数微分 \(d\ln Y/d\ln x\) が \(3\) より小さいあいだである。外側の井戸でその不等式が成り立つのは、\(\gamma_s(x)\) が \(2x\cosh x\sin x\) より大きいときとちょうど同じであり、両辺は \(1/2\) と \(3/5\) のあいだのただ一つの位相で等しくなる。その半径は \(5\ell/6\) と \(\ell\) のあいだに厳密にあり、第一の節の外にある。その半径から節までは、円軌道は不安定である。

円軌道のエネルギーは、向心の恒等式が定める運動エネルギーに、外側の井戸のポテンシャル~\eqref{eq:potential} を加えたものである。\(k e^{2}/\ell\) を単位にすると、それは \(x/\gamma_s(x)\) から、\(1/\gamma_s\) を零からその位相まで積分したものの二倍を引いた値である。戻る側では、位相が上がるにつれてエネルギーは下がる。すでに定めた一つの位相から節へ向かうと、エネルギーは上がる。その位相が、外側の井戸における円軌道のエネルギーの最小を担い、値は負である。そこでの力は、なお引力である。最小が定めるのは円軌道であり、円軌道を離れた落下の動径ポテンシャルは、節において下に有界でない。

核電荷を陽子の電荷で数えた数が陽子数 \(Z\) である。その核のまわりの電子一つは、陽子一つの力の \(Z\) 倍を受ける。向心の恒等式と外側の井戸のポテンシャルは、いずれも同じ因子を担う。したがって円軌道のエネルギーは、\(k e^{2}/\ell\) を単位にして、陽子一つのエネルギーの \(Z\) 倍であり、その半径での値は負のままである。力の対数微分は陽子一つの力の対数微分であるから、正の電荷はどれも、その半径の同じ側で円軌道を戻す。電荷を大きくすると、このエネルギーは低くなる。

\section{殻の電子数}
\label{sec:shells}

\(2\)、\(8\)、\(18\)、\(32\) は、等スケールの節の番号ではない。節はなお内側へ \(\Theta(1/n)\) と並ぶ。ボーアの半径の梯子 \(n^{2}\) は、そこからは読めない。

核電荷 \(Z\ge 1\) のまわりで、直径の両端に電子を二つ置く。外側の井戸に動径の釣り合いはない。半径も直径も第一の節と第二の節のあいだにある第一斥力殻には、\(Z\) がどれであっても、ちょうど一つある。そこでは核が外へ押し、二電子は引き合う。半径の大きい側では外向きの力は内を向き、半径の小さい側では外を向く。半径は戻る。角は戻らない。半径を固定すると、対の間隔は電子が対蹠点にあるとき最大であり、その間隔の干渉因子が負であるところでは、間隔とともに増すエネルギーはその点で極大になる。二電子を一つの席へ寄せる変位を、静電力は戻さない。

この欠如は直径に限らない。核までの距離と相互の間隔がすべて外側の井戸にある二電子は、核電荷が少なくとも \(1/4\) ならば、配置がどうであっても同時には力を消さない。向かい合う射線の上では、斥力は有限な半径のすべてで単位核電荷の引力の四分の一より厳密に小さく、どの核も少なくともその四分の一である。間隔の干渉因子が正であるところでは、対蹠点は、間隔が広がるにつれて下がるエネルギーの角の極小である。電子は広がる。その広がりは半径を戻さない。核が引く井戸での二電子の閉じた殻は、単極の席が二つであることである。符号つきの力の零ではない。

同じ核のまわりで、立方体の頂点に電子を八つ置く。半径と三つの間隔、すなわち辺、面の対角線、体の対角線が、すべて第一斥力殻に入る窓で、動径の釣り合いはやはり \(Z\ge 1\) ごとに一つある。外向きの力はその半径を戻す。

正四面体の四つの頂点に電子を置く。半径と辺がともに第一斥力殻に入る窓で、動径の釣り合いは核電荷が少なくとも \(1\) ならばちょうど一つある。正八面体の六つの頂点に電子を置く。半径と辺と向かい合う頂点がすべてその殻に入る窓で、動径の釣り合いはやはりそのような電荷ごとに一つある。外向きの力は、いずれの場合もその半径を戻す。四と六は頂点数であり、席の数ではない。六つの頂点は、双極の席の六つではない。

十八は、立方対称の一つの図形の頂点数ではない。そのような図形の頂点は六、八、十二である。十八は席の数である。単極は一つのモード、双極は三つ、トレースが零の四重極は五つ、トレースが零の八重極は七つである。双対スピノルとそのスピンを反転した相手が、各モードに席を一つずつ置く。次数ごとに見ると、その席は \(2\)、\(6\)、\(10\)、\(14\) である。次数 \(0\) までが \(2\)、次数 \(1\) までが \(8\)、次数 \(2\) までが \(18\)、次数 \(3\) までが \(32\) である。同じ算術は \(n=1,2,3,4\) に対する \(2n^{2}\) である。四つの主殻の電子数は、この四つの数である。この数は節の位置を定めず、十八点や三十二点の図形を力の極小だとはここでは主張しない。

席は、方向の上の模様であり、核のまわりの路ではない。単極は、どの方向でも同じ値をとる。双極は、その軸に垂直で核を通る平面を境に符号を変える。一方は正、他方は負であり、その平面は値が零になる方向の集まりである。四重極は、方向を反転しても値が変わらない。零でない四重極は、ある方向では正、別の方向では負であり、それらの張る平面の中の方向に、値が零になるものがある。八重極は、方向を反転すると符号を変える。逆方向の値は、もとの値の符号を反転したものである。模様を定めるのはモードである。

電子軌道と呼ばれるものは、電子配置である。席が電子を置き、その置き場所は軌道ではない。電子が動くことはある。外側の井戸には第\ref{sec:nocent}節の円軌道が残り、半径を戻す点を外れた変位はその半径をなお横切る。共有結合では位置が固定され、電子はその配置を占める。金属のイオンのあいだでは、電子にそのような位置はない（第\ref{sec:free}節）。量子力学も、電子軌道を定常な模様としてすでに読む。符号つきの力は、その読みを電子の位置そのものまで運ぶ。

直径が一つの席へ滑るのは、第\ref{sec:pauli}節の重なりである。重力の極がそれをはじき出す。符号つきの力が打ち消し、ヘッシアンが正になる軌跡は、二中心の正方形である。

\section{実験室の井戸}
\label{sec:lab}

原子の距離で散乱と分光学が測る符号は、外側の井戸の符号である。そこでは \(\gamma_s>0\) である。核は電子を引き、電子どうしは斥け合う。有限な半径から無限遠へ向かう路は節を横切らない。積分~\eqref{eq:potential} は、したがって有限な除去エネルギーを割り当てる。反転した層は逆である。同符号は引き合い、異符号は斥け合い、無限遠への路は極に断たれる。

第\ref{sec:shells}節の、半径を戻す点は、その反転した層にある。それらは電離の平衡ではない。立方体では、障害は幾何的であり、核電荷 \(Z\ge 3\) のすべてに対して成り立つ。三つの弦、すなわち辺、面の対角線、体の対角線は、いずれも半径より長い。電子間の位相は、したがって核の位相より小さい。区間 \((0,x_1)\) で干渉因子は狭義に減少するので、この三つの弦から組んだ応答は、遠方での値を厳密に下回る。その値は外向き係数の和
\[
\frac{3\sqrt{3}}{4}+\frac{3\sqrt{6}}{8}+\frac14
\]
であり、\(2\) と \(3\) のあいだに厳密にある。\(Z\ge 3\) では応答は電荷に出会わず、外向きの力は外側の井戸の全体で内を向いたままである。中性の立方体も、それより重い立方体も、実験室の符号が成り立つところでは閉じない。直径も、\(Z\ge 1\) のすべてについて、すでに同じである。係数は一つで、\(1/4\) である。この四分の一が、二電子の問題の全体である。どの配置でも、すべての距離が外側の井戸に留まるあいだ、核電荷が少なくとも \(1/4\) ならば二電子は同時には力を消さない。向かい合う射線の上では、斥力は有限な半径のすべてで単位核電荷の引力の四分の一を厳密に下回る。電子はなお広がる。間隔の干渉因子が正であるところでは対蹠点は角の極小であり、半径は戻らない。この井戸での二電子の殻は、第\ref{sec:shells}節の単極の席が二つであることである。

正四面体、正八面体、正二十面体も、同じ障害に会う。それらの弦はどれも半径より長い。したがって区間 \((0,x_1)\) では、干渉因子の比はいずれも \(1\) を厳密に下回り、応答は遠方での係数を厳密に下回る。正四面体の係数は \(3\sqrt{6}/8\) であり、\(1\) より小さい。正八面体の係数は \(\sqrt{2}+1/4\) であり、\(2\) より小さい。辺は四つで、スケールは \(1/\sqrt{2}\) であり、向かい合う頂点は一つで、スケールは \(1/2\) である。正二十面体の係数は \(\frac54\sqrt{5+2\sqrt{5}}+\frac14\) であり、\(5\) より小さい。各頂点から、近傍が五つ、さらに遠い頂点が五つ、向かい合う頂点が一つある。核電荷が少なくとも \(1\)、\(2\)、\(5\) ならば、応答は電荷に出会わず、外向きの力は外側の井戸の全体で内を向いたままである。正四面体、正八面体、立方体、正二十面体の、中性の殻は、実験室の符号が成り立つところでは閉じない。正四面体と正八面体は、直径や立方体と同じく、反転した層で一度閉じ、外向きの力はその半径を戻す。

外側の井戸で半径を戻す点は、同符号の弦のどれかが半径より短いときに現れる。図形を縮めると、その弦は半径がまだ節の外にあるうちに第一の節へ進み、対の斥力は外へ発散する。大きな間隔では、幾何的な係数が核電荷を下回るならば、同じ図形は内へ落ちる。この二つの領域のあいだで、応答は電荷を横切る。

五つの正多面体のうち、そのような弦をもつのは正十二面体だけである。各頂点から、三つの近傍が、半径の二乗の \(2(3-\sqrt{5})/3\) 倍の二乗距離にある。その弦のスケール \(R/d\) は \(1\) より大きい。二十の頂点の遠方での係数は、\(15/2\) と \(8\) のあいだに厳密にある。最も近い弦が第一の節に達するより前、応答は \(7\) を厳密に上回る。核電荷が \(7\) 以下ならば、どの対もなお外側の井戸にあるあいだ応答は電荷に出会わず、外向きの力はその区間の全体で外を向く。核電荷 \(Z\ge 8\) のすべてに対して、係数は電荷を下回るので、大きな間隔では力は内を向く。応答、すなわち弦の係数に、半径での \(\gamma_s\) を弦での \(\gamma_s\) で割ったものを掛けて足したものは、最も近い弦が第一の節に達するより前のちょうど一つの位相で \(Z\) に等しくなり、そこでの半径は \(7\ell/10\) と \(10\ell/7\) のあいだに厳密にある。電荷を大きくすると、この釣り合いはより小さい半径へ移る。より大きい半径では力は内を向き、より小さい半径では、どの対もなお外側の井戸にあるあいだ、外を向く。その根では動径方向の力は消え、無限遠への除去は積分~\eqref{eq:potential} である。

二十は頂点数であり、主殻ではない。\(2\)、\(8\)、\(18\)、\(32\) は、なお各モードの上の、双対スピノルとそのスピンを反転した相手の席である。八電子の殻を立方体として詰めると、閉じるのは反転した層だけである。正十二面体の根は、外側の井戸における二十の頂点の釣り合いである。四つの主殻の数の一つではない。

\section{電気陰性度}
\label{sec:eneg}

外側の井戸で、核の引力を正に数え、\(k e^{2}/\ell^{2}\) を単位にすると、その大きさは電荷に \(Y=4x^{2}/\gamma_s(x)\) を掛けたものである。位相が共通ならば \(Y\) は同じであり、しかも正であるから、電荷の大きい方がより強く引く。\(Y\) は位相について厳密に増加する。したがって、より小さい電荷がより大きい電荷を上回って引くのは、より大きい位相、すなわちより小さい半径に立つときだけである。

電荷 \(Z\) と \(W\) を長さ \(D\) の線分の両端に置き、電子をその線分の上に置く。線分全体の位相を \(\varphi=\ell/(2D)\) と書き、電子と電荷 \(Z\) のあいだにある線分の割合を \(t\) と書く。両方の距離が外側の井戸にあるのは、\(D\) が第一の節の半径の二倍より大きく、かつ \(t\) が \(\varphi/x_1\) と \(1-\varphi/x_1\) のあいだに厳密にあるときに限る。\(Z\) へ向かう力は \(Z\,Y(\varphi/t)-W\,Y(\varphi/(1-t))\) である。その区間で、この力は \(t\) について厳密に減少し、消える割合はちょうど一つである。

中点での力は \((Z-W)\,Y(2\varphi)\) である。\(Z>W>0\) ならば、そこでの電子は \(Z\) へ引かれ、打ち消しは \(W\) に厳密に近い。\(Z\) の側では力は \(Z\) を向き、\(W\) の側では力は \(W\) を向く。打ち消しは戻されず、席でもない。線分のうち、より大きい電荷へ引かれる部分の方が長い。電荷が等しいとき、打ち消しは中点であり、電子はより近い核へ引かれる。一端の電荷を増やすと、打ち消しは他端へ移る。

この力の向きが、より大きい電荷の電気陰性度である。実験室の尺度は読み取らない。遮蔽は導入せず、\(\ell\) は導出しない。

\section{排他としての重力斥力}
\label{sec:pauli}

重力チャネルは同じ干渉因子を用いる。源は電荷に結びついた対ではなく、慣性質量に結びついた運動学的な対である~\cite{dst-main}。二つの電子質量に対して、間隔を増す向きの力は
\begin{equation}
\label{eq:grav}
F_G=-\frac{G m_e^{2}}{\gamma_{s,G}(r)\,r^{2}}
\end{equation}
である。\(\gamma_{s,G}\) は質量チャネルの等スケール因子であり、長さはあらゆる質量の重力を支配する \(\lambda_\pi\) である。\(\gamma_{s,G}>0\) では質量は引き合う。\(\gamma_{s,G}<0\) では斥け合う。\(\gamma_{s,G}\) の最初の零点は、同じ導関数~\eqref{eq:gamma} によって単純である。内側から近づくと、\(F_G\) は単純極をもち、外を向く。

前置因子は小さい。CODATA の値は
\begin{equation}
\label{eq:ratioG}
\frac{G m_e^{2}}{k e^{2}}=2.401\times 10^{-43}
\end{equation}
を与える。単純極は、小さい前置因子によっては除かれない。第一の重力節からの内向きの分数変位を \(r=r_1(1-\varepsilon)\)、\(\varepsilon\downarrow 0\) と書く。このとき \(x-x_1=x_1\varepsilon+O(\varepsilon^{2})\) であり、
\[
\bigl|\gamma_{s,G}\bigr|
  =2x_1\cosh x_1\sin x_1\,\varepsilon+O(\varepsilon^{2})
\]
である。係数は \(\frac{\pi}{2\sqrt{2}}\cosh\frac{\pi}{4}\) と \(2\cosh 1\sin 1\) のあいだに厳密にある。数値では \(2.2260\) である。外向きの力が、同じ半径における非特異な電磁力のスケール \(k e^{2}/r^{2}\) を超えるのは、\(\varepsilon<1.078\times 10^{-43}\) のときである。電子質量が定めるのは、その近傍の幅である。近傍があるかどうかは定めない。

この比較は、電磁因子が同じ半径でそれ自身零ではないことを仮定する。二つのチャネルは別の位相対であるから、重力の節が電磁の節である必要はなく、次節の正方形は \(\lambda_\pi\) を選ばない。二つの因子がともに消えれば、それらの比は~\eqref{eq:ratioG} のままになり、電磁の片が支配する。ここで読む排他は、電荷チャネルが共有しない節における重力の極である。

両方の電子を一つの軌道の席に置くと、その間隔は節を横切らされる。力~\eqref{eq:grav} は外向きで非有界になるので、共有された席はただちに離れる。本文は、この斥力層がパウリ原理であると提案する。そこに入るフェルミオンは一つの席を占めることができず、量子統計への先立った訴えを要しない。共有できない席は、双対スピノルの繰り返された射線である。同じモードの第二の席は、スピンを反転した相手の直交する射線であり、二つを合わせたものが一重項である。その提案を、ここでは電子に採用する。スピン統計定理の導出ではなく、質量チャネルの長さは \(\lambda_\pi\) のままである。

\section{正方形と共有結合の力}
\label{sec:covalent}

電荷 \(+e\) の陽子を \(z=\pm R/2\) に置き、電荷 \(-e\) の電子を円筒半径 \(\rho\) に、中間面で互いに向かい合わせて置く。三つの間隔は、陽子間距離 \(R\)、電子間距離 \(r_{ee}=2\rho\)、そして四本が等しい電子--陽子距離
\[
r_{ep}=\sqrt{(R/2)^{2}+\rho^{2}}
\]
である。角運動量は割り当てない。電子は核のあいだで、一つの配置の固定した位置を占める。平衡とは、各陽子の軸力と各電子の円筒方向の力がともに消えることである。

\(z=+R/2\) にある陽子の軸力を \(R\) を増す向きに正と数え、一電子の円筒方向の力を \(\rho\) を増す向きに正と数えると、
\begin{align}
F_z
  &=k e^{2}\left(
    \frac{1}{\gamma_s(R)\,R^{2}}
    -\frac{R}{\gamma_s(r_{ep})\,r_{ep}^{3}}
  \right),
\label{eq:Faxial}\\
F_{\rho}
  &=k e^{2}\left(
    \frac{1}{\gamma_s(r_{ee})\,r_{ee}^{2}}
    -\frac{2\rho}{\gamma_s(r_{ep})\,r_{ep}^{3}}
  \right)
\label{eq:Frho}
\end{align}
である。\(F_{\rho}\) の第一項は電子間の力である。外側の井戸ではそれは斥力であり、外を向く。その外向きの片は静電的であり、遠心項が占めていた場所を占める。斥力殻では \(\gamma_s(r_{ee})<0\) であるから、同じ項は内を向く。そこでは同符号の電荷が引き合うからである。

四つの電荷が正方形の頂点であり、陽子が一方の対角線上、電子が他方の対角線上にある配置は
\begin{equation}
\label{eq:square}
\rho=\frac{R}{2},
\qquad
r_{ee}=R,
\qquad
r_{ep}=\frac{R}{\sqrt{2}}
\end{equation}
である。この軌跡の上で \(F_z\) と \(F_{\rho}\) は一つの式に還元される。\(x=\ell/(2R)\) と置けば
\begin{equation}
\label{eq:squareF}
F_z=F_{\rho}
  =\frac{k e^{2}}{R^{2}}
  \left(
    \frac{1}{\gamma_s(x)}
    -\frac{2\sqrt{2}}{\gamma_s(\sqrt{2}\,x)}
  \right)
\end{equation}
である。両方が消えるのは
\begin{equation}
\label{eq:squarecond}
\gamma_s(\sqrt{2}\,x)=2\sqrt{2}\,\gamma_s(x)
\end{equation}
のときである。係数 \(2\sqrt{2}\) は \(1\) より大きい。区間 \([0,\pi]\) の上で \(\gamma_s\) は減少する。対角線の位相 \(x\) がまだ第一の節より小さいあいだ、\(\gamma_s(x)\) は正であり、\(\sqrt{2}\,x\) は \(\pi\) に達していない。大きい方の位相での値は \(x\) での値より小さく、その値は自分自身の \(2\sqrt{2}\) 倍より小さい。釣り合いの両辺は出会わない。辺だけが節を越えていても、なお出会わない。一方の因子が負で、他方が正になるからである。対角線が外側の井戸にある正方形は、平衡ではない。

電子の間隔の、陽子の間隔に対する比を \(s\) と書く。正方形は \(s=1\) である。二つの間隔がともに外側の井戸にあるかぎり、二つの力は同時には消えない。同時に消えるためには、陽子の間隔での \(\gamma_s\) と電子の間隔での \(\gamma_s\) との比が \(s^{3}\) でなければならない。その井戸で \(\gamma_s\) は正であり、位相について狭義に減少する。\(s>1\) では電子の間隔の方が長く、その位相は小さいので、因子の比は \(1\) より厳密に小さい。\(s^{3}\) は \(1\) より厳密に大きい。\(s<1\) では二つの比較の向きが入れ替わる。\(s=1\) では因子の比は \(1\) であり、正方形の条件は厳密に負のままである。二つの因子の符号が反対のときも、その比は正の \(s^{3}\) と一致しない。三つの因子がどれも零でないとき、二つの力が同時に消えるのは、\(\gamma_s(R)\,R^{3}\)、\(\gamma_s(r_{ee})\,r_{ee}^{3}\)、\(\gamma_s(r_{ep})\,r_{ep}^{3}\) が一つの値であるときに限る。電子と陽子の間隔が外側の井戸にあるあいだ、そのような値はない。その井戸では、位相 \(\ell/(2r)\) の三乗で因子を割った商は、位相が上がるにつれて下がる。したがってその井戸にある二つの位相が同じ商をもつのは、それらが一致するときに限り、中央平面は残る位相を同じ井戸の別の値に置く。残る二つの位相がどちらも井戸にないならば、それぞれは少なくとも第一の節であり、電子と陽子の位相はその節の少なくとも \(\sqrt{2}\) 倍である。実験室の結合は、陽子が電子を引く配置である。その間隔は外側の井戸にあり、この対はそこにあるあいだ平衡をもたない。

釣り合いが成り立ちうるのは、対角線が第一の斥力殻に入ったあとである。そのとき対角線も辺も二つの節のあいだにある。そこでは両方の因子が負である。同符号の電荷は対角線に沿って引き合い、異符号の電荷は辺に沿って斥け合う。この範囲で
\[
\gamma_s(\sqrt{2}\,x)-2\sqrt{2}\,\gamma_s(x)
\]
は \(x\) とともに増す。\(3/2\) では負であり、\(8/5\) では正である。したがって殻の中で一度だけ、そして一度しか、零を通る。交点はその二つの位相のあいだにある。\(8/5\) での辺の位相はなお \(\pi\) より小さく、第二の節は \(\pi\) の向こうにあるので、交点は殻の内部にある。\(3/2\) と \(8/5\) のあいだでは両方の因子が負のままであり、正方形の上の共通の力は釣り合いと同じ符号をもつ。より長い対角線は内へ引かれ、より短い対角線は外へ押される。したがってこの根は正方形を戻す。その位相は
\[
x=1.55463,
\qquad
\frac{R}{\ell}=0.32162,
\qquad
\frac{r_{ep}}{\ell}=0.22742
\]
である。二つの干渉因子は、対角線上で \(\gamma_s(x)=-2.2208\)、辺上で \(\gamma_s(\sqrt{2}\,x)=-6.2812\) である。電子を陽子から離して支える斥力は、辺の力である。遠心力ではない。

この根における部分力は
\begin{align}
\bigl|F_{\mathrm{diag}}\bigr|
  &=\frac{k e^{2}}{\bigl|\gamma_s(x)\bigr|\,R^{2}}
  =4.3532\,\frac{k e^{2}}{\ell^{2}},
\label{eq:Fdiag}\\
\bigl|F_{\mathrm{side}}\bigr|
  &=\frac{k e^{2}}{\bigl|\gamma_s(\sqrt{2}\,x)\bigr|\,r_{ep}^{2}}
  =3.0782\,\frac{k e^{2}}{\ell^{2}}
\label{eq:Fside}
\end{align}
である。対角線の力が共有結合の力のスケールである。それは陽子間の間隔に沿う引力であり、辺の斥力の軸成分によって打ち消される。辺の力は、対角線の力を \(\sqrt{2}\) で割ったものである。

この根の近くで、静電ポテンシャルを \(R\) と \(\rho\) の関数とみる。
そのヘッシアンの固有値は \(37.33\,k e^{2}/\ell^{3}\) と \(7.933\,k e^{2}/\ell^{3}\) であり、どちらも正である。根は局所的な極小である。\(\rho\) が \(R\) に従うとき、対角線の剛性は
\begin{equation}
\label{eq:kappa}
\kappa=8.178\,\frac{k e^{2}}{\ell^{3}}
\end{equation}
である。\(F_z=F_{\rho}=0\) の根は、同じ殻の中にさらに二つある。それぞれヘッシアンの固有値の符号は反対であり、したがって鞍点である。正方形が、この族の安定な点である。その二つの電子は、一つの席ではなく、間隔 \(R\) の向かい合う頂点を占める。両方を一つの頂点に置く変位は、第\ref{sec:pauli}節の重なりである。その間隔が第一の重力節に達すると、\(F_G\) がそれらをはじき出す。電磁的なヘッシアンは、小さな変位に対して正方形を安定にする。重力の極は、\(m_e\) が小さくても、共有された席を禁ずる。根は殻の内部にあるので、外側の井戸の積分~\eqref{eq:potential} は、無限遠で分かれた原子に対する解離エネルギーをこれに割り当てない。

ヘリウムは電荷二の核と二つの電子であり、一中心の二電子の殻である。同じ中間面の族の各核に電荷二を置き、電子は二つに保つ。正方形の上では、軸方向の釣り合いは辺の干渉因子が対角線の因子の \(\sqrt{2}\) 倍であることを求め、円筒方向の釣り合いはその因子の \(4\sqrt{2}\) 倍であることを求める。したがって二つの因子が有限で零でないところでは、正方形は平衡ではない。核の間隔と電子の間隔がともに外側の井戸にあるあいだ、二つの力はどの比でも同時には消えない。電子の方が核より離れているところでは、因子の比は \(1\) より厳密に小さく、釣り合いは電子の間隔の核の間隔に対する比の三乗の四倍を求める。その三乗は \(1\) より厳密に大きい。電子の方が近いところでは、各電子と核との間隔は核どうしの間隔より短い。軸方向の釣り合いは、核の間隔の因子の \(\sqrt{2}\) 倍より大きい因子を、その短い間隔の上に求める。しかしその位相は第二の節に達しない。因子が正であるあいだは核の間隔の因子より小さく、位相が第一の節を越えると負になる。二つの間隔の因子の符号が反対のときも、釣り合いは正の比を求めるので成り立たない。実験室の符号は核が電子を引くことであり、この対はその井戸に平衡をもたない。電荷二と二つの電子は、一中心の二電子の殻のままである。単極の席が二つであることである。

\section{水の分子}
\label{sec:water}

水分子は、核を原点に置き、二本の射線のそれぞれに陽子を一つ、結合の電子を一つ置く。電子は、その陽子より核に近い。射線の外に、さらに電子が立ちうる。射線のあいだの角は変わりうる。一本の射線の上の電子と陽子は別の点である。

二つの結合電子が斥け合うあいだ、もう一方の結合は、各結合電子をより広い角へ押す。もう一方の射線の陽子は逆向きに引き、その引きは勝たない。この比較に核電荷は入らない。同じ射線の上の電子は、接線方向の力を加えない。結合電子を、射線の、もう一方の結合に向いた側から斥ける電子は、その電子をさらに広い角へ押す。反対側から斥ける電子は、より狭い角へ押す。さらに加わる電子のすべてが、その射線の上に立つか、もう一方の結合に向いた側から結合電子を斥けるならば、結合電子に働く接線方向の力は、より広い角の側に留まる。結合電子が陽子より核に近く、結合電子同士が斥け合うあいだ、その力は直線に達しない角のすべてで正であり、直線の上で消える。直線は結合電子を戻す。

各陽子に働く接線方向の力は、別の釣り合いである。核と、同じ射線の上の電子は寄与せず、核電荷は入らない。残るのは、もう一方の陽子の斥力から、もう一方の結合電子の引力を引いたものである。二つの間隔がともに外側の井戸にあり、陽子どうしの間隔が、各陽子ともう一方の結合電子との間隔を超えないあいだ、力はより広い角の側を向く。二つの間隔が等しくなる射線のあいだの角は、直線の三分の一より厳密に小さい。電子が核から陽子の少なくとも四分の一の距離にあり、その等間隔がなお外側の井戸にあるとき、同じ差は直線の上で負である。陽子に働く力は、等間隔ではより広い角の側を向き、直線に達しないより大きい角のどこかではより狭い角の側を向く。したがってその力は、等間隔と直線のあいだの角で消える。直線は陽子を戻さない。その角で結合電子の間隔がなお外側の井戸にあるならば、結合電子の接線方向の力はより広い角の側を向いたままであり、陽子に働く力が消える射線の上で結合電子は静止しない。

直線の上では、核からの二つの距離が外側の井戸にあるとき、二本の結合だけが及ぼす接線方向の力はない。二本の結合の反対側の軸に立つ電子は、その配置を直線から外す。結合電子が斥け合う直線に達しない角では、結合は各結合電子をより広い角へ押す。反対側の電子は、その間隔が外側の井戸にあるあいだ、より狭い角へ押す。直線の上で接線方向の力は厳密に負になり、直線は結合電子を戻さない。

その直線の上で、各電子を核とその陽子のあいだに置き、核からの二つの距離を外側の井戸に置き、電子とその陽子のあいだの区間の干渉因子を節でないとする。核電荷が少なくとも \(1/4\) ならば、電子に働く外向きの力と陽子に働く外向きの力の和は、厳密に負である。その区間に沿う引力は、和の中で打ち消される。水分子は、この二つの結合と、それらの上にない電子とである。

\section{共有結合の角}
\label{sec:angle}

核から出る各射線は、一つの電子対の電子を運ぶ。共有電子対の陽子は、同じ射線の、より外側にあり、その電子を外へ引く。非共有電子対の射線には陽子がない。正四面体の射線の上に四つの電子を一つの距離に置き、そのうちのいくつかの射線に陽子をより大きい共通の距離に置き、間隔のすべてを外側の井戸に置く。陽子のある射線の電子に働く外向きの力は、陽子のない射線の電子に働く外向きの力より厳密に大きい。非共有電子対の電子は内へ押され、より近い電子になる。

それらの射線のあいだの角の余弦は \(-1/3\) であり、角は直角より大きい。四つの電子が一つの距離にあり、四つの陽子がそれより大きい一つの距離にあるとき、接線方向の力はない。この釣り合った角が、四つの共有電子対の角である。角が直角以上であり、二つの間隔がともに外側の井戸にあるあいだ、電子を核に近づけると隣との間隔は縮み、位相は上がり、斥力 \(Y=4x^{2}/\gamma_s\) は強くなる。より近い電子が、隣の電子の核からの距離の少なくとも \(1/\sqrt{2}\) 倍の距離に留まるならば、その斥力の接線方向の部分も強くなる。

三つの電子を共通の距離に留める。四つ目の電子を、隣の核からの距離の \(1/\sqrt{2}\) 倍より内側へは入れずに内側へ移し、その射線の陽子を除く。間隔は外側の井戸に留める。余弦 \(-1/3\) では、残る三つの電子のそれぞれに働く接線方向の力は、内側へ移した電子から離れる向きを向き、三つのあいだの角は狭められる。もう一つの電子を同じだけ内側へ移し、その陽子も同じく除くと、狭める向きの寄与は同じになり、二つ合わせた寄与は一つの二倍である。核により近い非共有電子対は、結合距離にある共有電子対より強く、周りの結合を押しのける。そのような対が二つあるとき、押しのける力は一つの二倍であり、残る結合に働く力は、その角を四つの等しい対の角より下へ押す。射線が二本だけで、その外に何もないとき、第\ref{sec:water}節の直線は結合電子を戻す。水の二つの結合は、二本の共有の射線の外に、より近い電子が二つ立った場合である。

\section{配位と金属結合の力}
\label{sec:metal}

一価の配位殻で近傍が \(z\) 個あるとき、最近接の結合には電子含量 \(\nu=2/z\) を与える。二原子は \(z=1\)、\(\nu=2\) であり、第\ref{sec:covalent}節の正方形である。電荷 \(-\nu e\) を、中間面の向かい合う二点にある電荷 \(-(\nu/2)e\) の二つの塊に分ける。力~\eqref{eq:Faxial} と~\eqref{eq:Frho} は
\begin{align}
F_z
  &=k e^{2}\left(
    \frac{1}{\gamma_s(R)\,R^{2}}
    -\frac{\nu R}{2\,\gamma_s(r_{ep})\,r_{ep}^{3}}
  \right),
\label{eq:Fznu}\\
F_{\rho}
  &=k e^{2}\left(
    \frac{(\nu/2)^{2}}{\gamma_s(r_{ee})\,r_{ee}^{2}}
    -\frac{\nu\rho}{\gamma_s(r_{ep})\,r_{ep}^{3}}
  \right)
\end{align}
になる。平衡はなお \(F_z=F_{\rho}=0\) であり、遠心項はない。外側の井戸での外向きの片は \(F_{\rho}\) の第一項である。塊は \(\gamma_s(r_{ee})>0\) のところで斥け合う。

\(z=8\) と \(z=12\) では、この外側の井戸の方程式に根が一つずつある。三つの間隔はすべて \(r_1\) の外側にあるので、符号は通常のものである。同符号は斥け合い、異符号は引き合う。陽子間の部分力と、共有結合の対角線~\eqref{eq:Fdiag} に対するその比は
\begin{align}
z=8:\quad
\frac{R}{\ell}&=2.3591,&
F_{\parallel}&=0.18820\,\frac{k e^{2}}{\ell^{2}},&
\frac{F_{\parallel}}{|F_{\mathrm{diag}}|}&=0.04323,
\label{eq:ratio8}\\
z=12:\quad
\frac{R}{\ell}&=1.2666,&
F_{\parallel}&=0.74188\,\frac{k e^{2}}{\ell^{2}},&
\frac{F_{\parallel}}{|F_{\mathrm{diag}}|}&=0.1704
\label{eq:ratio12}
\end{align}
である。どちらの根でも、ポテンシャルのヘッシアンは反対符号の固有値をもつ。それぞれ鞍点であり、局所的な極小ではない。斥力殻の内部にあるさらなる根も同様である。電子含量を完全な対から \(2/z\) に減らすと、符号つきの力が打ち消し合う点は残る。安定な配置は残らない。鞍点は塊の模型の欠陥ではない。どのイオンに対しても外側の井戸にある点は、電子を保持しない（第\ref{sec:free}節）。経験的な最近接距離は、陽子間の項だけによって \(F_{\parallel}\) を定める。それは入力であり、比には要らない。

\section{自由電子}
\label{sec:free}

金属の電子は、どのイオンも電子を引くところ、すなわち各イオンの外側の井戸に立つ。そのような電子を、ある点からある方向へ変位させる。その方向に沿う力の成分は定まった割合で変わり、割合が正であることは、変位がさらに押し進められることを意味する。\(e\) を単位として数えた点電荷 \(Q\) が位相 \(x=\ell/(2r)\) にあるとき、互いに垂直な任意の三つの方向に沿う割合の和は
\begin{equation}
\label{eq:outflow}
Q\,\frac{k e^{2}}{\ell^{3}}\,S(x),
\qquad
S(x)=\frac{16x^{4}\cosh x\sin x}{\gamma_s(x)^{2}}
\end{equation}
である。外側の井戸では \(S\) は正であり、位相について厳密に増える。したがってイオンは正の和を寄与し、さらに加わる電子は負の和を寄与する。クーロン力では、同じ和は零である。その零がアーンショーの定理である。クーロンの法則のもとでは、どのような電荷の配置も、さらに加えた電荷をすべての方向に保持することはない。符号つきの因子はその零を除き、しかも電荷の二つの符号に対して反対の向きに除く。

したがって、どのイオンに対しても外側の井戸にある電子は、互いに垂直な任意の三つの方向のうち少なくとも一つに沿って、イオンによって変位をさらに押し進められる。さらに加わる電子は逆に働き、\(S\) は距離が縮むにつれて増す。さらに加わる電子のそれぞれが外側の井戸にあって別々のイオンより遠くにあり、各イオンの電荷が少なくとも一であるかぎり、イオンがなお勝つ。イオンのあいだのどの点も電子を保持しない。電子に働く力が打ち消し合うところでも、その打ち消しはすべての方向には戻されない。第\ref{sec:eneg}節は、線分に沿って同じことを見出した。

対称性は、割合をすべての方向で同じにする。正八面体、立方体、立方八面体に並んだイオンの中心では、そのうち一つか二つが欠けていてもよく、電子に働く力は消え、どの変位も一つの割合でさらに押し進められる。\(k e^{2}/\ell^{3}\) を単位として、各殻は、そのイオンの数 \(6\)、\(8\)、\(12\) の三分の一に、その電荷と、その位相での \(S\) とを掛けたものを寄与する。単純立方、体心立方、面心立方の配列の格子点のまわりの最も近い二つの殻は、これらの種類である。面心立方の配列の八面体の隙間のまわりの最も近い二つの殻、すなわち八面体とそれに続く立方体も、同じである。

同じ殻は、その中心に正のイオンを保持する。イオンに働く力は消え、イオンのどの変位も、電子の割合にイオンの電荷を掛けて符号を反転した割合で戻される。クーロンの法則のもとでは、どちらの割合も零である。外側の井戸では、イオンは同符号の近傍によって位置に保持され、電子はそのような位置のどれからも押し出される。実験室の井戸にある金属は、保持されたイオンの配列であり、電子はそのあいだに席をもたない。電子は自由である。第\ref{sec:covalent}節の共有結合の正方形は反対の場合である。その位置は反転した殻にあり、ヘッシアンは正定値である。

無限の格子和はとらず、電気伝導度も凝集エネルギーも導出しない。遮蔽は導入しない。各イオンは点電荷である。第\ref{sec:crystal}節は、立方格子の打ち切りのすべての殻を数える。

\section{結晶}
\label{sec:crystal}

\eqref{eq:outflow} の割合 \(S\) は、\(\sin x\) の符号をもつ。第一の節を除いて \((0,\pi)\) のどの位相でも正であり、その節のどちら側でもそうである。\(\pi\) で消え、\(\gamma_s\) が零でないかぎり \((\pi,2\pi)\) では負である。したがって、対称な格子点のまわりの同符号の電荷は、節のどちら側でも、半径 \(\ell/(2\pi)\) までは互いに保持し合い、位相が \(\pi\) を超えると互いに押し出す。反転するのはその半径であり、節ではない。\(\gamma_s=1\) では、三方向の和は零である。

格子点のまわりの点電荷の有限集合が、電荷ごと、正四面体の回転で自分自身に写されるならば、その格子点で力は消える。体対角軸のまわりの三回回転と、座標軸のまわりの半回転とで足りる。試験電荷 \(q\) を変位させると、外向き成分はすべての方向で一つの割合、すなわち \(-q/3\) に \(\sum Q\,S\) を掛けたもので変わる。単純立方格子、体心立方格子、面心立方格子の、格子点を中心とする立方体と任意の半径とで切った部分は、いずれもそのような集合である。どの電荷も正であり、どの位相も \(\pi\) より小さく、かつ第一の節の上にないあいだ、その格子点は正のイオンを保持し、電子を押し出す。位相は距離とともに小さくなる。したがって最近接のイオンが外側の井戸にあれば、より遠い殻もすべてそこにあり、保持はその部分の全体に及び、殻の数はいくらでもよい。無限和はとらない。

外側の井戸では、\(S(x)/x^{4}\) は厳密に増える。最も近い殻が異符号で、二乗距離 \(k_{1}\) に電荷 \(Q_{1}\) のイオンが \(N_{1}\) 個あり、次の殻が同符号で、二乗距離 \(k_{2}\) に電荷 \(Q_{2}\) のイオンが \(N_{2}\) 個あり、近い殻が外側の井戸にあるとする。四面体の対称性はなお力を消し、
\begin{equation}
\label{eq:twoshell}
N_{2}Q_{2}k_{1}^{2}\le N_{1}Q_{1}k_{2}^{2}
\end{equation}
ならば、正の試験電荷はすべての方向に押し出される。四つの配列の最も近い二つの殻はこの限界を満たし、ありふれた電荷も満たす。岩塩型は、二乗距離 \(1\) に異符号の近傍を六個、二乗距離 \(2\) に同符号の近傍を十二個もち、\eqref{eq:twoshell} は \(Q_{2}\le 2Q_{1}\) である。塩化セシウム型は二乗距離 \(3\) に八個、二乗距離 \(4\) に六個をもつ。閃亜鉛鉱型は二乗距離 \(3\) に四個、二乗距離 \(8\) に十二個をもつ。どちらの限界も \(27Q_{2}\le 64Q_{1}\) である。蛍石型を陽イオンから見ると、二乗距離 \(3\) に八個、二乗距離 \(8\) に十二個あり、限界は \(27Q_{2}\le 128Q_{1}\) である。陰イオンから見ると、二乗距離 \(3\) に四個、二乗距離 \(4\) に六個あり、限界は \(27Q_{2}\le 32Q_{1}\) である。等しい電荷は初めの三つの限界を満たす。陰イオンの二倍の電荷をもつ陽イオンは、蛍石型の両方の格子点を満たす。これは最も近い二つの殻だけである。同じ格子点でのクーロンの割合は零である。この二つの殻はイオンを保持しない。

格子を自分自身に写す線形写像の跡は整数である。角度 \(\theta\) の回転は、本義でも変義でも、跡 \(2\cos\theta+\varepsilon\)、\(\varepsilon=\pm 1\) をもつ。したがって \(2\pi/n\) の回転が三次元の格子を保つのは、\(n\in\{1,2,3,4,6\}\) のときに限る。結晶の格子点は五回軸をもたない。第\ref{sec:lab}節の正二十面体と正十二面体は、格子点の対称性ではない。一中心の殻としては、あの節の計算のままである。凝集エネルギーも電気伝導度も導出せず、各イオンは点電荷である。

\section{数値例}
\label{sec:numbers}

この節の長さは \(\ell\) の倍数であり、力は \(k e^{2}/\ell^{2}\) の倍数である。電磁の長さは、電子の換算コンプトン波長である。

実験室で一電子の半径に測られる符号は入力であり、この長さに上界を与える。ボーア半径 \(a_0\) での引力は、その半径を外側の井戸に置く。したがって \(\ell<2x_1 a_0\) である。第一の節は \(19/20\) を厳密に下回るので、上界は \(\ell<(19/10)a_0\) である。SI の代理値では、この量は \(1005\times 10^{-13}\,\mathrm{m}\) と \(1006\times 10^{-13}\,\mathrm{m}\) のあいだに厳密にある。核電荷 \(Z\) がなお電子を引くところでの水素様半径 \(a_0/Z\) では、同じ符号が上界を \(\ell<(19/10)a_0/Z\) まで締める。電荷を大きくすると上界は小さくなる。入力 \(Z=2\) では、上界は \(502\times 10^{-13}\,\mathrm{m}\) と \(503\times 10^{-13}\,\mathrm{m}\) のあいだに厳密にある。入力 \(Z=92\) では、\(1092\times 10^{-15}\,\mathrm{m}\) と \(1093\times 10^{-15}\,\mathrm{m}\) のあいだに厳密にある。第一の節を \(a_0\) に置くと、その半径で力は発散し、半径 \(a_0/2\) の位相は \(\gamma_s<0\) のところに来る。\(a_0/2\) での引力は、その配置を外側の井戸の外に置く。

電磁の長さは、電子の換算コンプトン波長 \(\hbar/(m_e c)\) である。SI の代理値では、これは \(3861\times 10^{-16}\,\mathrm{m}\) と \(3862\times 10^{-16}\,\mathrm{m}\) のあいだに厳密にあり、電荷 \(92\) の上界を厳密に下回る。その第一の節は、この長さの \(10/19\) と \(15/28\) のあいだに厳密にあり、したがって \(2032\times 10^{-16}\,\mathrm{m}\) と \(2069\times 10^{-16}\,\mathrm{m}\) のあいだに厳密にあり、\(a_0/92\) の内側に厳密にある。半径 \(a_0\)、\(a_0/2\)、\(a_0/92\) は外側の井戸に留まる。上界は、この波長の下にある。

円軌道の半径、第一の節、対角線を戻す正方形の対角線、その辺、一電子の極値、第二の節は、この順序に並び、それぞれ次のものより外にある。第\ref{sec:nocent}節の円軌道の比較が変わる位相は \(1/2\) と \(3/5\) のあいだに厳密にあり、その半径は \(5\ell/6\) と \(\ell\) のあいだに厳密にある。第一の節は \(14/15\) と \(19/20\) のあいだに厳密にあり、その半径は \(10\ell/19\) と \(15\ell/28\) のあいだに厳密にあり、\(15\ell/28\) は \(5\ell/6\) より厳密に小さい。第\ref{sec:covalent}節の正方形は \(153/100\) と \(157/100\) のあいだに厳密にあり、対角線は \(50\ell/157\) と \(50\ell/153\) のあいだに厳密にあり、\(50\ell/153\) は \(10\ell/19\) より厳密に小さい。辺は対角線を \(\sqrt{2}\) で割った長さであり、\(50\ell/157\) を厳密に下回り、\(\ell/5\) を厳密に上回る。極値は \(\ell/(2\pi)\) と \(\ell/5\) のあいだに厳密にある。第二の節は \((\pi+\pi/4,\pi+\pi/2)\) にあり、その半径は \(\ell/(3\pi)\) と \(2\ell/(5\pi)\) のあいだに厳密にあり、\(2\ell/(5\pi)\) は \(\ell/(2\pi)\) より厳密に小さい。値 \(x_1=0.93755\)、\(r_1/\ell=0.53330\)、位相 \(1.55463\) と対角線 \(0.32162\,\ell\)、および \(r_2/\ell=0.12731\) は、これらの区間にある。

正方形の対角線の力は \(4x^{2}/\bigl|\gamma_s(x)\bigr|\) であり、区間 \((153/100,157/100)\) の上では \(4\) より厳密に大きい。\(x=\pi\)、すなわち \(\gamma_s=-\cosh\pi\) では、異符号の対に働く外向きの力は \(4\pi^{2}/\cosh\pi\) であり、\(3\) より厳密に大きく、\(7/2\) より厳密に小さい。共有結合の対角線は、したがって \(\gamma_s\) の極値における一電子の力を上回る。辺の力は対角線の力を \(\sqrt{2}\) で割ったものであり、\(2\sqrt{2}\) より厳密に大きい。値 \(4.3532\)、\(3.4057\)、\(3.0782\) は、これらの限界を満たす。

原子の距離で測られる符号は引力であるから、ボーア半径は外側の井戸にあり、そこでは陽子が電子を引く。正方形は反転した層にあり、辺に沿っては陽子が電子を斥ける。その辺をボーア半径に置くと、第一の節は \(\sqrt{2}\,x/x_1\) ボーア半径の位置に来て、比は \(91/40\) と \(12/5\) のあいだに厳密にある。実験室の符号は、この同定を許さない。解離エネルギーは掲げない。平衡は斥力殻の中にあり、\eqref{eq:potential} は節を越えない。

\section{符号の境界を横切る振動}
\label{sec:heat}

半径を戻す点は、内向きの力と外向きの力との境界である。半径の大きい側では外向きの力が内を向き、半径の小さい側では外を向く。根から外れた変位は、その符号の変化を横切って振動する。節は、その振動の席ではない。節では力が定義されず、以下の路は \(u\) が有限であるところに留まる。

対称な殻の一つの電子に働く外向きの力は
\begin{equation}
\label{eq:Fout}
F_{\mathrm{out}}(r)=\frac{k e^{2}}{r^{2}}\,u(x),
\qquad
x=\frac{\ell}{2r}
\end{equation}
である。\(u\) は、直径と正十二面体にすでに用いた、核の項と弦の項の組み合わせである。単純な根 \(u(x_*)=0\) で \(u'(x_*)>0\) ならば、
\begin{equation}
\label{eq:curv}
F_{\mathrm{out}}'(r_*)=-\frac{k e^{2}}{\ell^{3}}\,K,
\qquad
K=8x_*^4\,u'(x_*)
\end{equation}
である。電子質量は、この曲率を振動数に変える。
\begin{equation}
\label{eq:omega}
\omega^{2}=\frac{K\,k e^{2}}{m_e\ell^{3}}.
\end{equation}
この線形の読みでは、\(\omega\) は \(K\) と \(m_e\) と \(\ell\) とで定まる。振幅とともに増すことはなく、温度とともに増すこともない。

動径の根は、角速度を自由に残す。極座標では
\begin{equation}
\label{eq:polar}
m_e\bigl(\ddot r-r\dot\theta^{2}\bigr)=F_{\mathrm{out}}(r)
\end{equation}
である。項 \(m_e r\dot\theta^{2}\) は、接線速度の運動学である。\eqref{eq:force} への新しい寄与ではない。速度が小さいあいだ、根は外へ
\begin{equation}
\label{eq:shift}
\delta r=\frac{m_e\langle v_{\theta}^{2}\rangle}{|F_{\mathrm{out}}'(r_*)|\,r_*}
\end{equation}
だけずれる。接線方向のエネルギーが動径振幅 \(A\) のエネルギーに等しいとき、ずれは \(A^{2}/r_*\) であり、\(A\) より因子 \(A/r_*\) だけ小さい。平均の半径は、静電的な根のままである。電子は殻を周回し、そのあいだ半径は符号の境界を横切って振動する。直径では、第\ref{sec:shells}節のとおり、角は極大である。その周回それ自体は振動ではない。振動は、動径が境界を横切る運動である。

温度は、この振動子のエネルギーである。動径の一つのモードに対する古典的な等分配は
\begin{equation}
\label{eq:amp}
\langle(\delta r)^{2}\rangle
  =\frac{k_B T}{|F_{\mathrm{out}}'(r_*)|}
  =\frac{k_B T\,\ell^{3}}{K\,k e^{2}}
\end{equation}
を与える。二乗平均平方根の振幅は \(\sqrt{T}\) とともに増す。その増大が、運動の激しさである。定数 \(k_B\) と、以下に現れる \(c\) および \(h\) は、実験室の定数である。\(\gamma_s\) はそれらを定めない。

加速された電荷は放射する。非相対論的な電子一つのラーモア電力の平均は
\begin{equation}
\label{eq:larmor}
P=\frac{2k e^{2}\langle a^{2}\rangle}{3c^{3}},
\qquad
a=-\omega^{2}\,\delta r
\end{equation}
である。\eqref{eq:amp} により \(\langle a^{2}\rangle=\omega^{2} k_B T/m_e\) であるから、\(P\) は \(\omega\) を固定すれば \(T\) に比例し、\(T\) を固定すれば \(K\) に比例する。同じ根をより熱く占有すれば、同じ振動数でより多くの電力を放射する。熱を持つ物体は、そのような根の大きな集合である。その熱線は、この放射である。一つの根の輝線ではない。

すでに手元にある曲率は、互いに遠く離れている。電荷 \(8\) の正十二面体は、第\ref{sec:lab}節のただ一つの二十頂点の根であり、\(K=4.556\) をもつ。電荷 \(1\) の直径は、反転した層の中にあり、\(K=126.45\) をもつ。正方形の上で陽子を固定すると、\(\rho\) 方向の曲率は \(18.10\,k e^{2}/\ell^{3}\) であり、\(K\) の位置にこの数を置いて~\eqref{eq:omega} が使える。対角線の剛性~\eqref{eq:kappa} を担うのは、陽子の換算質量 \(m_p/2\) である。

古典的な電力は、より硬い根を優遇する。\(P\) が \(K\) とともに増すからである。熱を持つ物体の熱線は、以下に読む空洞のスペクトルであり、間隔が \(k_B T\) に比べて大きければ振動子は空になる。実験室の変位則は、\(300\,\mathrm{K}\) の物体のピークを \(9.659\,\mu\mathrm{m}\) に置く。そこでは \(k_B T=0.02585\,\mathrm{eV}\) である~\cite{codata2018}。

ラーモアの公式に入る振動数は、温度が割り当てられるより前に、符号の境界によって固定される。動く座標の慣性を \(m_e\) の位置に置けば、\eqref{eq:omega} がそのすべてである。根の無次元曲率、その慣性、そして \(\ell\) である。温度は、根を占有するエネルギーである。それは振幅~\eqref{eq:amp} を定め、\(a=-\omega^{2}\,\delta r\) を通じて電力~\eqref{eq:larmor} を定める。占有が振動数を組み直すことはない。

一つの符号つきの釣り合いが、なお二つの振動数で放射しうる。その上を二つの座標が動きうるからである。正方形の上で、陽子を固定した電子の円筒運動は、曲率 \(18.10\,k e^{2}/\ell^{3}\) と慣性 \(m_e\) とをもつ。対角線は、剛性 \(8.178\,k e^{2}/\ell^{3}\) と慣性 \(m_p/2\) とをもつ。\eqref{eq:omega} により、波長の比は、電子の座標における \(K/m\) と対角線における \(K/m\) との比の平方根である。その比は \(45.08\) である。どちらの波長も \(\ell^{3/2}\) に比例するので、この比は辺の同定を用いない。電子が核の周期で加速されるのは、辺が対角線に拘束されているからであり、符号の境界を電子自身が横切る曲率は \(18.10\,k e^{2}/\ell^{3}\) のままである。

温度が動かすのは、これらの根がすでに作っている順序の上の占有である。その順序は \(K/m\) の順序である。\(300\,\mathrm{K}\) では、変位のピークは \(0.1283\,\mathrm{eV}\) にあり、これは \(4.96\,k_B T\) である。温度を上げると \(k_B T\) が上がり、符号つきの力の、より高い根が占有に入りうる。通常の温度に属する激しさは、この順序の遅い端、軟らかい境界である。線そのものは固定されたままである。連続ピークの変位は、間隔を \(k_B T\) の単位で測る梯子の占有である。

遷移は、その振動数の内部にこのような運動をもたない。条件
\begin{equation}
\label{eq:bohrnu}
\nu=\frac{E_i-E_j}{h}
\end{equation}
は、\(\nu\) を一対の固有値から取る。固有値はハミルトニアンのスペクトルである。温度は、そのスペクトルの上のアンサンブルを径数づけ、スペクトルは変えない。したがって差は、冷たい物体でも熱い物体でも同じである。変わるのは速さであり、それは上の状態の占有を通じてである。\(K/m\) と \(\ell\) が定める差は、その意味で不変である。熱線のピークは別の対象である。それは \(4.96\,k_B T\) に座り、温度に比例する。ピークを遷移と同定することは、一つの振動数に、固有値の差であることと \(T\) に比例することとを同時に要求することである。固定されたハミルトニアンのスペクトルは、そのような比例を許さない。これが矛盾である。

熱線は、熱を持つ物体の黒体放射である。定常軌道は電子の運動を固定する。陽子が運ぶのは、原子の熱速度である。二つの速度の差は、軌道が計算する量ではなく、温度とともに変わらない。黒体の波長は、温度とともに変わる。軌道の描像は、したがって黒体を説明しない。

そのような線の一つの占有は、等比級数である。間隔を \(x=\varepsilon/(k_B T)\)、\(x>0\) と書く。基底状態より上の量子の平均個数は
\[
\langle n\rangle=\frac{1}{e^{x}-1}
\]
であり、基底状態より上の平均エネルギーは \(\varepsilon\langle n\rangle\) である。\(x\) が零に近づくと、このエネルギーは \(k_B T\) に戻る。\(x\ge 1\) では、量子の平均個数は \(2e^{-x}\) より小さい。間隔が小さいところでは、温度を固定した古典的な電力はなお \(K\) とともに増え、より硬い根ほど多く放射する。間隔が大きいところでは、同じ重みが振動子を空にする。間隔そのものは、符号の境界が定めた振動数~\eqref{eq:omega} のままである。

大きな空洞は、振動数を共有するモードの数で重みをつけた、そのような線の和である。空洞では、区間 \(d\nu\) に入るモードの数が \(\nu^{2}\,d\nu\) に比例して増える。この増え方は空洞の幾何である。プランクの \(h\) は、振動数を間隔 \(x=h\nu/k_B T\) に変える外部の定数として残る。\(k_B\) と \(c\) も外部のままである。\(d\nu\) の中のエネルギーは、したがって \(x^{3}/(e^{x}-1)\) に比例する。この形は正の最大を一つもち、それは \((3-x)e^{x}=3\) のただ一つの正の解であって、\(5/2\) と \(3\) のあいだにある。最大の振動数は温度に比例する。

波長で数えると、\(\nu=c/\lambda\) は因子 \(c/\lambda^{2}\) を持ち込む。形は \(x^{5}/(e^{x}-1)\) になり、\(x=hc/(\lambda k_B T)\) である。そのただ一つの最大は \(9/2\) と \(5\) のあいだにあり、\((5-x)e^{x}=5\) のただ一つの正の解である。実験室のピーク \(4.96\,k_B T\) は、この波長の最大である。振動数の最大は、より小さい倍数であり、\(3\,k_B T\) より下にある。最大を定める \(x\) の値は \(T\) によらない。したがってピークの波長と温度との積は定数である。温度が変えるのは、一つの固定された梯子を読む単位である。この定数の積が、ウィーンの変位則である。

\(T\nu^{2}\) に比例する密度、すなわちレイリー--ジーンズの法則は、どの振動数でも増加するので、空洞の上の等分配は最大をもたない。最大が現れるのは、等比級数の重み \(1/(e^{x}-1)\) が分配 \(k_B T\) の位置に立ったときである。

振動数の形の積分は
\[
\int_{0}^{\infty}\frac{x^{3}}{e^{x}-1}\,dx=\frac{\pi^{4}}{15}
\]
である。空洞の、積分した振動数密度は、したがって \(T^{4}\) に比例する。一つの古典的な根の放射電力は \(T\) に比例する。この四乗は、その積分についてのシュテファン--ボルツマンの法則である。それを流束に変える因子、なかでも空洞の二つの偏波は、この勘定の外に残る。熱を持つ物体の熱線は、この積分である。半径を戻す根は、固定された振動数の順序のままであり、それぞれ~\eqref{eq:omega} を保つ。通常の温度が占有するのは、なおその順序の重い端である。衝突は、符号の境界上の一つの場所から別の場所へ電子を移しうる。そこで放射される場は、その根のラーモア場である。重み \(1/(e^{x}-1)\) が決めるのは、順序のどこまで寄与が生き残るかである。連続スペクトルを担うのは空洞である。

\section{結論}

\(\gamma_s\) の符号は、\eqref{eq:force} における斥力である。遠心項は必要でなく、導入もしない。電子一つと陽子一つには、二つの干渉因子の符号が同じところでは、力の消える半径がない。ボーア半径では、質量の項は電荷の項より、\(4.40\times 10^{-40}\) と \(4.42\times 10^{-40}\) のあいだに厳密にある因子だけ小さい。外側の井戸では、円軌道は \(5\ell/6\) と \(\ell\) のあいだに厳密にある半径より外で動径方向に戻り、その半径と第一の節のあいだでは不安定である。その井戸の円軌道のエネルギーは、同じ半径で最も小さく、値は負である。正の核電荷をもつ核のまわりの電子一つは、その半径を保ち、軌道のエネルギーは陽子一つの値にその電荷を掛けたものである。一つの核のまわりの直径、正四面体、正八面体、立方体は、第一の斥力殻に半径を戻す釣り合いを一つずつもち、直径は角の極大であり、\(18\) と \(32\) は四重極と八重極までのモードに双対スピノルとそのスピンを反転した相手の席を一つずつ置いた数であって節ではない。席は方向の上の模様である。電子軌道と呼ばれるものは、軌道ではなく電子配置である。電子が動くことはあり、共有結合では位置が固定される。単極だけが方向によらず、それより高い席は、ある方向では正、別の方向では負である。外側の井戸にある二電子は、核電荷が少なくとも \(1/4\) ならば同時には力を消さない。向かい合う射線の上では斥力は単位核電荷の引力の四分の一より厳密に小さく、間隔の干渉因子が正であるところでは対蹠点は角の極小である。その井戸での二電子の殻は、単極の席が二つであることであり、符号つきの力の零ではない。正四面体、正八面体、立方体、正二十面体は、核電荷が少なくとも \(1\)、\(2\)、\(3\)、\(5\) ならば、外側の井戸に根をもたない。正十二面体はもつ。その遠方での係数は \(15/2\) と \(8\) のあいだに厳密にあり、最も近い弦が第一の節に達するより前、応答は \(7\) を厳密に上回る。核電荷が \(7\) 以下ならば、その区間の全体で外へ押される。\(Z\ge 8\) のすべてに対して、その応答は、最も近い弦が第一の節に達するより前のちょうど一つの位相で電荷と等しくなり、半径は \(7\ell/10\) と \(10\ell/7\) のあいだに厳密にあり、電荷を大きくするとこの釣り合いはより小さい半径へ移り、無限遠への除去は外側の井戸に留まる。第一の節の半径の二倍より長い線分では、電子は中点からより大きい電荷の方へ引かれ、打ち消しは一つで、より小さい電荷に近く、戻されない。それぞれ二つずつあるときには、第一の斥力殻の中の正方形がちょうど一つあり、その正方形は対角線を戻す。その位相は \(3/2\) と \(8/5\) のあいだにあり、対角線の力は \(4.3532\,k e^{2}/\ell^{2}\)、剛性は \(8.178\,k e^{2}/\ell^{3}\) である。陽子の間隔と電子の間隔がともに外側の井戸にあるあいだは、その二つの間隔のどの比でも、二つの力は同時には消えない。電子と陽子の間隔が外側の井戸にあるあいだも、二つの力は同時には消えない。ヘリウムは電荷二の核と二つの電子であって二電子の殻のままであり、各核に電荷二を置いた同じ族は外側の井戸に平衡をもたず、二つの干渉因子が零でないかぎり正方形も平衡ではない。第\ref{sec:water}節の水分子も、結合電子が陽子より核に近く、斥け合うあいだ、さらに加わる電子のすべてが射線の上に立つか、もう一方の結合に向いた側から斥けるかぎり、接線方向の力をより広い角の側に留め、直線はそれらの電子を戻す。直線の上では、核電荷が少なくとも \(1/4\) ならば外向きの力の和は負である。電子が核から陽子の少なくとも四分の一の距離にあり、陽子の等間隔がなお外側の井戸にあるとき、各陽子に働く接線方向の力は、その等間隔と直線のあいだで消え、等間隔における射線のあいだの角は直線の三分の一より厳密に小さい。正四面体の射線の上の四つの共有電子対は、余弦 \(-1/3\) の角で接線方向の力を及ぼさない。その角は直角より大きい。非共有電子対の電子は、その射線に陽子がないために内へ押され、隣の核からの距離の少なくとも \(1/\sqrt{2}\) 倍に留まるあいだ、結合距離にある共有電子対より強く残る結合を押しのける。そのような電子が二つあるときの狭める寄与は一つの二倍である。電子対の \(2/8\) と \(2/12\) を運ぶ結合には静電的な根があり、ヘッシアンは不定である。それらの部分力は共有結合の対角線の \(0.04323\) と \(0.1704\) である。どのイオンに対しても外側の井戸にあるところでは、イオンだけでは電子を保持しない。互いに垂直な三つの変位がさらに押し進められる割合の和は正であり、クーロンの法則ではそれは零である。別々のイオンより遠くにある電子が加わっても、これは覆らない。正八面体、立方体、立方八面体に並んだイオンの中心では、電子はすべての方向に押し出され、正のイオンはすべての方向に戻される。イオンは保持され、電子は自由である。同符号のイオンからなる単純立方、体心立方、面心立方の立方体による打ち切りは、どの位相も \(\pi\) より小さく第一の節の上にないかぎり、殻をいくつ数えてもイオンを保持し電子を押し出し、割合は半径 \(\ell/(2\pi)\) で反転する。岩塩型、塩化セシウム型、閃亜鉛鉱型、および蛍石型の両方の格子点では、最も近い二つの殻がありふれた電荷に対して外側の井戸で正のイオンを押し出し、結晶の格子点は五回軸をもたない。

正方形は局所的な極小であり、その電子は向かい合う頂点に座る。重力チャネルが補う安定性は次である。二つの電子質量を第一の重力節の内側へ進めると、それらは単純極に出会い、共有された席ははじき出される。因子 \(G m_e^{2}/(k e^{2})=2.401\times 10^{-43}\) はその近傍を薄くし、近傍を消さない。節はなお、外側の井戸の積分が解離エネルギーを割り当てることを妨げる。ボーア・スペクトルは含意されない。正方形の辺はボーア半径ではない。その辺をその半径に置くと、第一の節は \(91/40\) と \(12/5\) ボーア半径のあいだに厳密に来る。その半径での実験室の符号は引力であるから、その半径は外側の井戸にある。その符号は \(\ell<(19/10)a_0\) を強制し、\(a_0/Z\) での同じ符号は \(\ell<(19/10)a_0/Z\) を強制する。入力の電荷 \(1\)、\(2\)、\(92\) に対して、上界はそれぞれ \(1005\times 10^{-13}\) と \(1006\times 10^{-13}\) メートル、\(502\times 10^{-13}\) と \(503\times 10^{-13}\) メートル、\(1092\times 10^{-15}\) と \(1093\times 10^{-15}\) メートルのあいだに厳密にある。第一の節を \(a_0\) に置くと、半径 \(a_0/2\) は \(\gamma_s<0\) のところに来る。電磁の長さは \(\hbar/(m_e c)\) であり、電荷 \(92\) の上界の下に厳密にあり、\(3861\times 10^{-16}\) と \(3862\times 10^{-16}\) メートルのあいだにあり、その第一の節は \(2032\times 10^{-16}\) と \(2069\times 10^{-16}\) メートルのあいだにあり、\(a_0/92\) の内側に厳密にある。正方形は反転した層にある。円軌道の半径、第一の節、対角線、辺、一電子の極値、第二の節は、この順序に並び、それぞれ次のものより外にある。

半径を戻す点は符号の境界である。それを横切る動径の振動は \(\omega^{2}=K k e^{2}/(m_e\ell^{3})\) をもち、電荷 \(8\) の正十二面体では \(K=4.556\) である。正方形の上では、電子の波長と対角線の波長との比は \(45.08\) である。電子の座標における \(K/m\) と対角線における \(K/m\) との比の平方根であり、辺の同定をこの比は要しない。一つの根の熱線は、その根がすでに持つ周期における、電荷の加速である。間隔が \(x\ge 1\) ならば平均の占有は \(2e^{-x}\) より小さい。遷移の振動数は固有値の差であり、温度に比例しない。定常軌道は電子の運動を固定する。陽子に対する電子の速度の差は、軌道が計算する量ではなく、温度とともに変わらない。連続スペクトルは、モードが振動数の二乗に比例して増える空洞の上で、占有の重みをつけた積分である。その波長の密度は、\(k_B T\) を単位として \(9/2\) と \(5\) のあいだに最大を一つもつ。\(300\,\mathrm{K}\) の物体のピークは \(0.1283\,\mathrm{eV}\) にあり、それはその最大の実験室の値である。ピークの波長と温度との積は定数である。積分した振動数密度は \(T^{4}\) に比例し、一つの古典的な根は \(T\) に比例する。モードの数と \(h\) は入力である。

\begin{thebibliography}{9}

\bibitem{dst-main}
Anonymous,
``Gravity as Torsion between Dual Spacetime: A Biquaternionic Reformulation of General Relativity'' (2025).
\url{https://github.com/hypernumbernet/dual-spacetime-doc/blob/main/double-spacetime-theory.pdf}.

\bibitem{codata2018}
E.~Tiesinga, P.~J.~Mohr, D.~B.~Newell, and B.~N.~Taylor,
``CODATA recommended values of the fundamental physical constants: 2018,''
\textit{Rev.\ Mod.\ Phys.} \textbf{93}, 025010 (2021).

\end{thebibliography}

\end{document}
